\documentclass[11pt]{article}

\usepackage[margin=1in]{geometry}
\usepackage{amsmath}
\usepackage{booktabs}
\usepackage{array}
\usepackage{siunitx}
\usepackage{microtype}
\usepackage{graphicx}
\usepackage{hyperref}
\hypersetup{colorlinks=true, linkcolor=blue, urlcolor=blue, citecolor=blue}

\newcolumntype{R}{>{\raggedleft\arraybackslash}p{1.6cm}}

\title{Which Parts of the ``Modern Transformer Recipe'' Actually Matter?\\
A Controlled Ablation on TinyStories}
\author{Cohen-Transformer}
\date{June 2026}

\begin{document}
\maketitle

\begin{abstract}
Contemporary language-model training stacks bundle a long list of small
architectural and optimization tweaks --- Muon, QK-norm, learnable QK-gain,
LayerScale, value residuals, z-loss, logit softcapping --- and present them as a
single ``modern recipe.'' We disentangle this bundle with a controlled,
multi-seed ablation. A fixed ${\sim}97.6$M-parameter decoder-only transformer is
trained on TinyStories under an identical token budget (30.72M tokens, 1000
steps) across 22 configurations $\times$ 3 seeds $=$ 66 runs, structured as a
fractional factorial design (all-off and all-on anchors, single-feature add-ins,
leave-one-out drops from the full recipe, and targeted $2\times2$ interaction
cells). Two components dominate: switching AdamW $\to$ Muon and adding QK-norm
each cut final validation loss by ${\approx}0.13$ nats, and the two are strongly
\emph{synergistic} --- together they buy ${\approx}0.24$ nats, far more than the
sum of their individual effects. LayerScale and value residuals each contribute
a smaller but real ${\approx}0.05$ nats and are mildly synergistic with each
other. QK-gain and logit softcap are effectively free riders ($|\Delta| < 0.002$
nats), and z-loss is mildly harmful at this scale ($+0.016$ nats). The best
single configuration is in fact the full Muon recipe \emph{minus} z-loss (1.651
val loss), not the kitchen-sink ``everything on'' (1.669). All 66 runs trained
to completion with zero divergences.
\end{abstract}

\section{Introduction}

The recipes shared in open LLM-training projects are cumulative: each new trick
is added on top of the last, and the resulting stack is evaluated as a whole.
This makes it hard to know which ingredients carry the gains, which are inert,
and which are actively counterproductive at a given scale. It also makes the
recipes hard to port --- a practitioner copying the full stack inherits its dead
weight and its failure modes.

This paper takes the seven toggles in the \texttt{Cohen-Transformer} training
stack and measures each one in isolation and in combination, holding everything
else fixed (model size, data, token budget, learning-rate schedule, weight
decay, gradient clipping, batch order). The question is deliberately narrow and
answerable: \emph{under a fixed compute budget, what does each component
contribute to validation loss, and do the components interact?}

\section{The components under study}

All seven are independent toggles over a common backbone (12 layers, 12 heads,
$d_{\text{model}}=768$, SwiGLU MLP, RMSNorm pre-norm, RoPE, untied LM head). The
backbone with all toggles off is the AdamW baseline.

\begin{table}[h]
\centering
\small
\begin{tabular}{@{}lp{8.2cm}l@{}}
\toprule
\textbf{Toggle} & \textbf{What it does} & \textbf{Where} \\
\midrule
optimizer & Muon (Newton--Schulz-orthogonalized momentum, \texttt{match\_rms\_adamw} LR) on 2D body matrices + AdamW on embeddings/head/scales, vs.\ pure AdamW & optimizer \\
qk\_norm & Per-head RMSNorm on Q and K before RoPE (Gemma-2 style) & attention \\
qk\_gain & Learnable per-head multiplicative gain folded into Q (scales scores pre-softmax) & attention \\
layerscale & Learnable per-channel residual gates on the attention and MLP branches (CaiT), init 0.1 & residual \\
value\_residual & Each layer mixes its value with layer-0's value via a learnable per-layer $\lambda$ (ResFormer, arXiv:2410.17897) & attention \\
z\_loss & \texttt{lse\_square\_scale}~$=1\mathrm{e}{-4}$ regularizer on the log-partition & loss \\
softcap & tanh logit softcap at 30.0 & loss \\
\bottomrule
\end{tabular}
\end{table}

z-loss and softcap are loss-side options of
\texttt{LigerFusedLinearCrossEntropyLoss}; the rest are model/optimizer changes.
The Muon $+$ AdamW split and weight-decay policy match the production
\texttt{simple.py} stack.

\section{Experimental setup}

\paragraph{Model.} ${\sim}97.6$M parameters (the small variation across
configurations --- 97{,}536{,}768 to 97{,}573{,}787 --- is the handful of extra
scale/gate parameters some toggles add; it is ${<}0.04\%$ and does not confound
the comparison).

\paragraph{Data.} \texttt{roneneldan/TinyStories}, tokenized once with a fixed
8192-token byte-level BPE into a memory-mapped \texttt{uint16} corpus. Every run
reads identical $(x,y)$ blocks in identical order (no shuffling, deterministic
batches), so the only differences between runs are the toggles and the seed.

\paragraph{Budget.} 1000 optimizer steps, batch size 15, gradient accumulation
1, block size 2048 $\to$ 30{,}720 tokens/step $\to$ \textbf{30.72M tokens per
run}. LR peak $1\mathrm{e}{-3}$, 99-step linear warmup then cosine decay to
$0.1\times$ peak, weight decay 0.1, gradient-norm clip 1.0, bf16 autocast,
\texttt{torch.compile}. Validation loss is the mean cross-entropy over 100
held-out blocks at the final step.

\paragraph{Design.} A full $2^7$ grid (128 configs) was not run. Instead a
fractional factorial focuses the budget on the questions that matter:
\begin{itemize}
  \item \textbf{Anchors} --- all-off (pure AdamW baseline) and all-on (full Muon recipe).
  \item \textbf{Add-one} --- baseline + exactly one toggle, for clean main effects of the features expected to matter on their own.
  \item \textbf{Leave-one-out} --- full recipe with exactly one toggle removed, to measure each component's \emph{marginal} contribution inside the complete stack.
  \item \textbf{Targeted $2\times2$ cells} --- Muon$\times$QK-norm, LayerScale$\times$value-residual, z-loss$\times$softcap, z-loss$\times$QK-gain, softcap$\times$QK-gain, for interaction terms.
\end{itemize}
22 unique configs $\times$ 3 seeds $=$ 66 runs. Main effects are computed over
matched on/off pairs (configs identical except the factor in question);
interactions over matched $2\times2$ quadruples. Effects are reported as
$\Delta(\text{final val loss}) = \text{ON} - \text{OFF}$, so \textbf{negative
means the component helps}.

\section{Results}

\subsection{Main effects}

\begin{table}[h]
\centering
\small
\begin{tabular}{@{}lRRRl@{}}
\toprule
\textbf{Factor} & \textbf{mean $\Delta$} & \textbf{std} & \textbf{n pairs} & \textbf{verdict} \\
\midrule
\textbf{qk\_norm}        & $-0.128$ & 0.070 & 9  & large help \\
\textbf{optimizer (Muon)}& $-0.128$ & 0.079 & 9  & large help \\
value\_residual          & $-0.053$ & 0.032 & 9  & moderate help \\
layerscale               & $-0.051$ & 0.034 & 9  & moderate help \\
qk\_gain                 & $-0.002$ & 0.008 & 15 & negligible \\
softcap                  & $-0.000$ & 0.010 & 15 & negligible \\
\textbf{z\_loss}         & $+0.016$ & 0.008 & 15 & mildly harmful \\
\bottomrule
\end{tabular}
\end{table}

Two findings stand out. First, QK-norm and the Muon optimizer are the workhorses
of the recipe, each worth ${\approx}0.13$ nats --- an order of magnitude more
than the next tier. Second, \textbf{z-loss consistently hurts} at this scale:
its effect is positive in mean and tight in variance (std 0.008 over 15 pairs),
so this is a real signal, not noise. z-loss is a numerical-stability regularizer
designed for large-vocab, long-schedule training; with an 8192-token vocab and a
1000-step budget there is no instability to suppress (zero runs diverged), and
the penalty simply pulls the model off the cross-entropy objective.

\subsection{Interactions}

$\text{interaction} = [\text{effect of A} \mid \text{B on}] - [\text{effect of A}
\mid \text{B off}]$; \textbf{negative $=$ synergy} (A helps more when B is
present), positive $=$ redundancy/antagonism.

\begin{table}[h]
\centering
\small
\begin{tabular}{@{}llRl@{}}
\toprule
\textbf{A} & \textbf{B} & \textbf{interaction} & \textbf{reading} \\
\midrule
optimizer  & qk\_norm        & $-0.158$ & strong synergy \\
layerscale & value\_residual & $-0.071$ & synergy \\
softcap    & qk\_gain        & $-0.007$ & {\raise.17ex\hbox{$\scriptstyle\sim$}}none \\
z\_loss    & softcap         & $+0.006$ & {\raise.17ex\hbox{$\scriptstyle\sim$}}none \\
z\_loss    & qk\_gain        & $+0.003$ & {\raise.17ex\hbox{$\scriptstyle\sim$}}none \\
\bottomrule
\end{tabular}
\end{table}

The Muon$\times$QK-norm interaction is the headline. The $2\times2$ cell means
make it concrete:

\begin{table}[h]
\centering
\small
\begin{tabular}{@{}lRR@{}}
\toprule
 & \textbf{QK-norm off} & \textbf{QK-norm on} \\
\midrule
\textbf{AdamW} & 1.987 & 1.929 \\
\textbf{Muon}  & 1.964 & \textbf{1.749} \\
\bottomrule
\end{tabular}
\end{table}

Alone, Muon buys ${\approx}0.023$ nats and QK-norm ${\approx}0.058$ nats;
together they buy ${\approx}0.238$ --- roughly three times the sum of the parts.
Muon orthogonalizes the body weight updates, and QK-norm bounds the scale of the
query/key projections feeding attention; the data say these two stabilizations
are complementary, and that \textbf{most of the ``modern recipe'' advantage is
concentrated in this single pair.} LayerScale and value residuals show the same
pattern more weakly: each gates or re-routes the residual stream, and they
reinforce each other ($-0.071$).

\subsection{Run rankings and the kitchen-sink trap}

The best individual configurations are all Muon $+$ QK-norm based:

\begin{table}[h]
\centering
\small
\begin{tabular}{@{}clR@{}}
\toprule
\textbf{rank} & \textbf{configuration} & \textbf{final val loss} \\
\midrule
1  & Muon, all on \textbf{except z-loss}      & \textbf{1.651} (seed mean 1.653) \\
4  & Muon, full recipe (all on)               & 1.669 (seed mean 1.670) \\
19 & Muon $+$ QK-norm only                    & 1.747 \\
22 & Muon, all on except QK-norm              & 1.781 \\
-- & AdamW baseline (all off)                 & 1.987 \\
-- & AdamW full recipe                        & 1.847 \\
\bottomrule
\end{tabular}
\end{table}

The top three runs (all three seeds) are the full recipe \emph{with z-loss
removed}, beating the kitchen-sink configuration by ${\approx}0.018$ nats ---
consistent with the main-effect sign of z-loss. Dropping QK-norm from the full
recipe (rank 22, 1.781) is the single most damaging leave-one-out, again
confirming QK-norm as the load-bearing component. Note also that the full recipe
under \textbf{AdamW} (1.847) never reaches even the \emph{worst} Muon recipe run,
underscoring how much of the gain flows through the optimizer$\times$QK-norm
channel.

End to end, the recipe moves the model from 1.987 (baseline) to 1.651 (best), a
0.336-nat reduction, of which the Muon$\times$QK-norm pair accounts for the large
majority.

\begin{figure}[h]
\centering
\includegraphics[width=\linewidth]{DataOutput/analysis/loss_curves_paper.png}
\caption{Validation-loss curves, one line per configuration (mean over 3 seeds,
with a shaded min--max band; the bands are narrow because seed spread, std
${\approx}0.003$--$0.008$ nats, is small next to the between-config gaps) for the
two anchors --- the all-off AdamW baseline (dashed black) and the all-on full
Muon recipe (green) --- and the leave-one-out runs (full recipe with one
component removed, labelled by the dropped component). The full recipe separates
from the baseline early and holds the gap throughout; among the leave-one-outs,
dropping Muon or QK-norm is by far the most damaging, while dropping z-loss
slightly \emph{improves} on the full recipe.}
\label{fig:curves}
\end{figure}

\section{Discussion}

\paragraph{A practitioner's shortlist.} If you must port only part of this
stack, port Muon and QK-norm \emph{together} --- separately they underwhelm,
together they deliver. LayerScale and value residuals are a worthwhile second
tier. QK-gain and softcap can be dropped with no measurable cost at this scale,
and z-loss should be dropped outright unless a longer schedule or larger vocab
reintroduces the instability it is meant to control.

\paragraph{Why z-loss flips sign.} z-loss and softcap both exist to tame logit
blow-up over long training. This study's regime (short budget, small vocab, zero
divergences) is exactly the one where that insurance is unnecessary, so it shows
up as pure cost (z-loss) or pure no-op (softcap). The result is a reminder that
``stability'' components should be judged against an actually-unstable baseline,
not folded into a recipe by default.

\paragraph{Scope and limitations.}
\begin{itemize}
  \item \textbf{Scale.} One model size (${\sim}98$M), one corpus (TinyStories --- simple, low-entropy text), one short budget (30.72M tokens). The \emph{signs} of the large effects are likely robust, but the inert/harmful verdicts for z-loss and softcap are the components most likely to flip at larger vocab or longer schedules --- precisely their design regime. These are the priority for a follow-up.
  \item \textbf{Fractional design.} Three- and higher-way interactions are not estimated; main effects and the five targeted pairs are. The strong Muon$\times$QK-norm synergy means single-factor numbers should be read as context-dependent, not additive.
  \item \textbf{Single metric.} Final validation loss only --- no downstream-task or generation-quality evaluation, and no measurement of wall-clock cost (Muon and the extra parameters add a small per-step overhead, visible in the logged \texttt{wall\_time\_s} but not analyzed here).
  \item \textbf{Seeds.} Three seeds give a usable variance estimate; the per-factor std columns show the large effects sit well outside seed noise, while QK-gain and softcap sit inside it.
\end{itemize}

\section{Conclusion}

The ``modern transformer recipe'' is not monolithic. On a controlled TinyStories
ablation, its benefit concentrates almost entirely in \textbf{Muon $+$ QK-norm
acting together}, with a smaller LayerScale $+$ value-residual contribution and a
trailing set of components that are inert (QK-gain, softcap) or mildly
counterproductive (z-loss) at this scale. The single best configuration is the
full recipe with z-loss removed. The broader lesson is methodological: stack
ingredients earn their place against matched, multi-seed pairs and explicit
interaction tests --- not by accumulation.

\appendix
\section{Reproduction}

\begin{verbatim}
python DataScripts/prepare_data.py          # one-time tokenization
python DataScripts/run_matrix.py --seeds 3  # 66 runs, resumable
python DataScripts/analyze_results.py       # tables + figures
\end{verbatim}

Per-run artifacts: \texttt{DataOutput/runs/<run\_id>/\{config.json, log.jsonl,
summary.json\}}. Aggregated tables and loss-curve figures:
\texttt{DataOutput/analysis/\{results.csv, main\_effects.csv, interactions.csv,
summary.md, loss\_curves*.png\}}. All 66 runs completed; 0 diverged.

\end{document}
